\documentclass[10pt]{article}

\usepackage[utf8]{inputenc}

\usepackage[T1]{fontenc}

\usepackage{amsmath}

\usepackage{amsfonts}

\usepackage{amssymb}

\usepackage[version=4]{mhchem}

\usepackage{stmaryrd}



\begin{document}

множестве, снабженном какой-либо другой структурой. Так, говорят об естественной топологии метрич. пространства или об естественной (интервальной) топологии упорядоченного множества. Первая имеет своей базой множество всех открытых паров данного метрич. пространства, вторая - открытые интервалы данного линейно упорядоченного множества.



Т. п. наз. м е т р и з у е м ы м, если во множестве $X$ eго точек можно ввести метрику $\rho$, порождающую топологию данного Т. П. Метризуемые пространства образуют один из важнейших классов Т. П., и к центральным проблемам общей топологии принадлежали в течение нескольких десятилетий общая и специальная проблемы метризации, т. е. проблемы нахождения необходимых и достаточных условий для того, чтобы Т. П. или Т. П. того или иного специального класса было метризуемым. Эти условия составляют содержание общих или специальных метризационных теорем.



Всякое подмножество $X_{0}$ множества $X$ всех точек данного Т. П. $X$, Ж естественно превращается в Т. п. (подпространство пространства $X$, (\$) с топологией, әлементы к-рой суть всевозможные множества вида $X_{0} \cap \Gamma$, где $\Gamma$ - любой элемент топологии (5). Пусть дана (открытая) топология (5) в множестве $X$, превращающая это множество в Т. П. X, (\$. Т. К. топология есть множество нек-рых подмножеств множества $X$, то между различными топологиями в одном и том же множестве $X$ естественно устанавливается отношение (частичного) порядка (по включению), т. е. топология $\mathfrak{H}_{2}$ больше (или равна) топологии $\mathfrak{H}_{1}$, если $\mathfrak{H}_{1}$ есть подмножество множества $\mathfrak{S H}_{2}$, т. е. каждое множество, открытое в топологии $\$_{1}$, будет открытым и в топологии (\$). Если в пределах данного рассуждения речь идет об одной и той же топологии в данном множестве $X$, то обыкновенно Т. П. $X$, (S) обозначается просто $X$. Из понятия топологии выводятся и все остальные основные топологич. понятия. Прежде всего замкнутые множества определяются как дополнения $\mathrm{k}$ открытым. Далее, о к р е с т н о с т ь ю точки $x$ в данном пространстве $X$ наз. всякое открытое множество, содержащее точку $x$. Понятие окрестности позволяет сразу же определить и понятие т о ч к и п р и к о с н ов е н и я для любого множества $M \subset X$ как такой точки, любая окрестность к-рой имеет непустое пересечение с множеством $M$. Из этого определения следует, что любая точка самого множества $M$ является точкой прикосновения этого множества. Множество всех точек прикосновения множества $M$ наз. з а м ы к а н и е м множества $M$ и обозначается $[M]$. Переход от любого множества $M$ к его замыканию наз. о п е р а ц и е й з а мы к а н и я в данном Т. П. Свойства этой операции: 1) $M \subset[M]$, причем $M=[M]$ тогда и только тогда, когда $M$ замкнуто, т. е. его дополнение открыто; 2) $\left.\left[M_{1} \cup M_{2}\right]=\left[M_{1}\right] \cup\left[M_{2}\right] ; 3\right)[[M]]=[M]$. Замыкание любого множества есть пересечение всех замкнутых множеств, содержащих множество $M$; другими словами, $[M]$ есть наименьшее замкнутое множество, содержащее множество $M$.



Операция замыкания и ее основные свойства 1), 2), 3) получены, исходя из основного (в нашем изложении) понятия топологии или топологич. структуры, введенной в данное множество $X$. Можно было бы, наоборот, считать основным топологич. понятием понятие замыкания, т. е. считать, что в абстрактно данном множестве $X$ для каждого подмножества $M$ определено подмножество $[M]$, наз. замыканием множества, так что выполнены свойства 1), 2), 3) (наз. в этом случае а к си о ма ми 3 а мык ания, или а к сиом ам и К у р а т о в с ко г о) и 4) $\varnothing=[\varnothing]$. На основе так введенного понятия замыкания замкнутые множества определятся как множества, совпадающие со своими замыканиями, а открытые множества - как множества, дополнительные к замкнутым; таким образом получается в точности топология в нашем первоначальном смысле, а олерация замыкания, к к-рой она приводит, совпадает с той, к-рая была дана априори. Именно этот путь и был избран К. Куратовским (К. Киratowski, 1922) для построения понятия Т. п. В 1925 П. С. Александровым были введены открытые топологич. структуры. Оба подхода приводят к одному и тому же классу Т. П., в настоящее время являющемуся общепринятым.



С понятием Т. п. тесно связано понятие непрерывного отображения одного пространства в другое. Отсбражение $f: X \rightarrow Y$ Т. п. $X$ в Т. п. $Y$ н е п р е р ы в но в точк е $x \in X$, если для любой окрестности $O_{y}$ точки $y=f(x) \in \bar{Y}$ в пространстве $Y$ существует такая окрестность $O_{x}$ точки $x$ в $X$, что $f\left(O_{x}\right) \subset O_{y}$ (у с л о в и е К о ші и). При этом, не изменяя содержания определения, можно брать окрестности $O_{y}$ и $O_{x}$ из любых открытых баз соответственно пространств $Y$ и $X$. В частности, для метрич. пространств это определение непрерывности переходит в обычное известное из курсов математич. анализа определение. Если отображение $f: X \rightarrow Y$ непрерывно в каждой точке $x \in X$, то оно наз. н е п ре ры вным о тоб р а ж ени ем пространства $X$ в пространство $Y$. Для непрерывности отображения $f: X \rightarrow Y$ каждое из следующих условий необходимо и достаточно. 1) Если $x$ есть точка прикосновения какого-либо множества $M \subset X$, то $f(x)$ есть точка прикосновения множества $f(M)$ в $Y$. 2) Полный прообраз $f^{-1}(\Gamma)$ всякого открытого в $Y$ множества $\Gamma$ есть открытое множество в $X$. Аналогично для замкнутых множеств .



Если дано (непрерывное) отображение $f$ Т. п. $X$ в Т. П. $Y$ и $X_{0}$ есть подпространство пространства $X$, то в силу отображения $f$ пространство $X_{0}$ отображается в $Y$ и это отображение (наз. о г р а н и ч е н и е м о тоб р а же ни я $f: X \rightarrow Y$ на подпространство $\left.X_{0}\right)$ также непрерывно.



Важный частный случай непрерывных отображений образуют т. н. Ф а к т о р ны е о т о б р а ж е н и я, характеризующиеся каждым из следующих эквивалентных между собой условий. Множество $B \subset Y$ открыто (замкнуто) в $Y$ тогда и только тогда, когда $f^{-1}(B)$ обладает тем же свойством в $X$. Если непрерывное отображение $f$ пространства $X$ на пространство $Y$ взаимно однозначно, то определено и обратное отображение $f^{-1}: Y \rightarrow X$, но это отображение может не быть непрерывным. Если же оно непрерывно, то каждое из отображений $f, f^{-1}$ взаимно однозначно отображает топологию одного из пространств $X, Y$ на топологию другого пространства. Тогда каждое из двух взаимно однозначных отображений $f$ и $f^{-3}$ наз. т о п о л о г и ч ес ким отображ е и е м, или гоме омо рф и з м о м. Два пространства $X$ и $Y$, из к-рых одно может быть гомеоморфно отображено на другое, наз. гом е о м о рыным и или топологи ч е с ки ә к в и в а л е н т ны м и между собою.



Непрерывное отображение $f: X \rightarrow Y$ наз. н е п р ив о д и м ы м, если при этом отображении никакое отличное от всего $X$ замкнутое множество $M$ не отображается на все $Y$.



Конкретное изучение Т. П. связано прежде всего с выделением из общего класса этих пространств подклассов, характеризующихся теми или иными дополнительными условиями или аксиомами, кроме первоначальных аксиом, определяющих топологию пространства. Эти дополнительные аксиомы бывают разной природы. Прежде всего это группа т. н. а к с и о м о т д е л и м о с т и. Первой аксиомой отделимости была а к с и о м а X а у с д о р ф а, заключающаяся в требовании, чтобы любые две различные точки про-





\end{document}